% ===============================================
% ECEV 33400: Stochastic Processes      Fall 2017
% hw_draft.tex
% Template for first draft of homework submission
% ===============================================


% -----------------------------------------------
% The preamble that follows can be ignored. Go on
% down to the section that says "START HERE" 
% -----------------------------------------------

\documentclass{article}

\usepackage[margin=1in]{geometry} 
\usepackage{amsmath,amsthm,amssymb,enumitem}
\usepackage[sectionbib]{natbib}
\usepackage{chemformula}
\usepackage{mathtools}

\newcommand{\R}{\mathbb{R}}
\newcommand{\Z}{\mathbb{Z}}
\newcommand{\N}{\mathbb{N}}
\newcommand{\Q}{\mathbb{Q}}
\newcommand{\C}{\mathbb{C}}
\newcommand{\F}{{\mathbb F}}
\newcommand{\Pb}{{\mathbb P}}
\newcommand{\E}{{\mathbb E}}
\newcommand{\Frac}[2]{\left(\frac{#1}{#2}\right)}
\newcommand{\Eval}[2]{\big|_{#1}^{#2}}
\newcommand{\intinf}{\int_{-\infty}^{\infty}}

\newenvironment{theorem}[2][Theorem]{\begin{trivlist}
\item[\hskip \labelsep {\bfseries #1}\hskip \labelsep {\bfseries #2.}]}{\end{trivlist}}
\newenvironment{lemma}[2][Lemma]{\begin{trivlist}
\item[\hskip \labelsep {\bfseries #1}\hskip \labelsep {\bfseries #2.}]}{\end{trivlist}}
\newenvironment{exercise}[2][Exercise]{\begin{trivlist}
\item[\hskip \labelsep {\bfseries #1}\hskip \labelsep {\bfseries #2.}]}{\end{trivlist}}
\newenvironment{problem}[2][Problem]{\begin{trivlist}
\item[\hskip \labelsep {\bfseries #1}\hskip \labelsep {\bfseries #2.}]}{\end{trivlist}}
\newenvironment{question}[2][Question]{\begin{trivlist}
\item[\hskip \labelsep {\bfseries #1}\hskip \labelsep {\bfseries #2.}]}{\end{trivlist}}
\newenvironment{corollary}[2][Corollary]{\begin{trivlist}
\item[\hskip \labelsep {\bfseries #1}\hskip \labelsep {\bfseries #2.}]}{\end{trivlist}}

\newenvironment{solution}{\begin{proof}[Solution]}{\end{proof}}

\begin{document}

% ------------------------------------------ %
%                 START HERE                 %
% ------------------------------------------ %

\title{Assignment 4} % Replace X with the appropriate number
\author{ \iffalse STUDENT NAME \fi \\ECEV 33400} % Replace  "\iffalse X \fi" with your name.
\date{Due Date: Thursday 11/02/2017} % Don't change this! :( 

\maketitle

% -----------------------------------------------------
% The following two environments (theorem, proof) are
% where you will enter the statement and proof of your
% first problem for this assignment.
%
% In the theorem environment, you can replace the word
% "theorem" in the \begin and \end commands with
% "exercise", "problem", "lemma", etc., depending on
% what you are submitting. Replace the "x.yz" with the
% appropriate number for your problem.
%
% If your problem does not involve a formal proof, you
% can change the word "proof" in the \begin and \end
% commands with "solution".
% -----------------------------------------------------

% \subsection*{Time Dependence and the Probability Generating Function}

% \begin{enumerate}
% \item 
% \end{enumerate}

\subsection*{Michaelis-Menten kinetics and the Holling Type II Functional Response}
\begin{enumerate}
\item Consider a tiny reactor of volume $V$ containing a solution of chemical species $A$ and $E$. $E$ is an enzyme catalyzing the conversion of $A$ into a product $W$ through the following catalytic, non-reversible reactions:
%%% \ch[math-space=0pt]{ E + S <=>[ $k_1$ ][ $k_{-1}$ ] ES <=> EP <=> E + P }
\begin{equation}
A \,+ \,E  \xrightarrow{k_1} AE \xrightarrow{k_2} E \,+\, W \nonumber
\end{equation}    
However, chemical species $W$ has a tendency to immediately precipitate and disappear from the solution. Consider that the concentration of the enzyme, in each of its two chemical forms, $E$ or $EA$ is constant, $[E] + [EA] = \Phi_E$. The concentration of the chemical species $A$ at the initial time is $[A]= \Phi_A(0) = \Phi_0$. 
\begin{itemize}
\item Write first the system of ODEs describing the temporal evolution of the concentration of each of the chemical species involved in the reaction, the {\em macroscopic law}. 
\item Describe the state of the system through the number of molecules of species A, $n$, and the number of molecules of free enzyme, $E$, $m$. Assume that the total number of molecules of enzyme either free or forming the substrate-enzyme molecule is constant, $M$.   
\item Write a two dimensional master equation for the temporal evolution of the probability of having $n$ and $m$ molecules of species $A$ and enzyme, respectively. 
\item Write the total rate for the system undergoing any configuration change. 
\item Consider that there are $n_0$ molecules of species $A$ and $M$ molecules of free enzyme at the initial time. Take $k_1 = 1$, $k_2 = 10$, $n_0 = 10000$, and $M=100$, and $V=100$. Build stochastic simulations and check that averages of each of the chemical species over stochastic simulations matches its deterministic evolution as described by the {\em macroscopic law}.  
\item After the initial transient, calculate the temporal evolution for the average of the ratio of the number of free enzyme molecules to the number of substrate-enzyme molecules over stochastic simulations. Is this ratio reaching, approximate, a stationary value? What is the best prediction for this average at this quasi-stationary state?  
\end{itemize}
\end{enumerate}
\bibliography{ecology_david}
\bibliographystyle{unsrt}
% -----------------------------------------------------
% First problem
% -----------------------------------------------------

% \begin{problem}{1}
% \end{problem}

% \begin{solution}
% \end{solution}

% -----------------------------------------------------
% Second problem
% -----------------------------------------------------

% \vspace{0.25in} % This just adds some space between problems 1 and 2.

% \begin{problem}{2}
% \end{problem}

% \begin{solution}
% \end{solution}

% -----------------------------------------------------
% Third problem
% -----------------------------------------------------


% -----------------------------------------------------
% Fourth problem
% -----------------------------------------------------


% -----------------------------------------------
% Ignore everything that appears below this.
% -----------------------------------------------

\end{document}
